\section{Discussion}
\label{sec:discussion}

\paragraph{What closes the studio-to-field gap.} For an ImageNet CNN trained on studio leaves, two
simple properties of the pooling account for the whole architectural gain we observed: knowing where
the leaf is, and pooling locally within it. Each is worth 5--8 points on raw field photographs, they
add up to 12--14 points, and both survive without rotation augmentation and on top of strong
background-compositing augmentation. Background compositing and mask pooling act similarly (each lifts
plain ConvNeXt-T to about 41--42\%), which supports the reading that the first part is about discounting
the background. The second part, localisation, helps on raw photographs and on texture corruptions but
hardly on whitened photographs, so it seems to make features robust to clutter and texture loss inside
and around the leaf rather than to add shape information; shape bias barely changes.

\paragraph{What does not.} We expected measured leaf geometry to matter: an object-centred sweep,
per-slab area and boundary length of the blade and venation, and attention restricted to the slab
being measured. None of these ingredients improves transfer beyond localised pooling. Object-centred
slabs are no better than image-axis strips, the same measurements fused globally give nothing, and
restricting attention to the measured slab can be replaced by global attention. Within \SVA{}, a fine
sweep and learned weights matter, but only relative to weaker versions of \SVA{} itself. Euler
tokens make a model sensitive to holes in the blade, yet transfer worse. Within the limits of our
design, the effect we can detect is about \emph{where} features are pooled, not about measuring the
leaf. We report \SVA{} because a well-controlled negative result for a natural and increasingly
popular idea, fusing topological and geometric measurements with deep features, is useful, and
because its localised variant is among the best models we trained.

\paragraph{Practical recommendations.} (i)~Evaluate leaf classifiers on field photographs; studio
accuracy is saturated and uninformative (98--100\% for every model). (ii)~With an ImageNet CNN, pool
over the leaf and locally within it, e.g.\ with strip pooling, which is simpler and cheaper than \SVA{}
(\cref{app:cost}), and combine it with background-compositing augmentation; together they raised
field accuracy from 34\% to 53--55\%. (iii)~Where a biology foundation model such as BioCLIP~2 can be
used, use it: it reaches 81\% without any module, and none of the modules changes it.

\paragraph{Limitations.} (i)~One studio dataset and one field dataset with 20 shared species. Field
labels come from vernacular folder names that we mapped to species and checked visually; residual
label noise would lower all models alike. The 20-way accuracy uses the knowledge of which species can
occur. (ii)~The masked field condition and all localised models depend on our masker, whose accuracy
on field photographs we could not measure without ground-truth masks; the raw-photo results, which the
plain model obtains without any mask, are the primary outcome. (iii)~ConvNeXt-T results use five
seeds, EfficientNet-B0 and the key ablations three, ViT-S and BioCLIP~2 one; differences of 2--3 points
between localised models are not resolved by our design, and the ViT results are indicative only.
(iv)~One fixed training recipe was used for all models of a backbone; a recipe tuned per model could
shift individual numbers. (v)~The venation field is learned from ridge-filter pseudo-labels and is not
validated against traced veins.
